\documentclass[14pt]{extarticle}
% =============================================================================
% Shared Preamble for Worksheets, Homeworks, Exams & Lesson Notes
% (scripts/tex_preamble.tex)
% =============================================================================

\newif\ifdetailed
\ifdefined\setdetailed
  \setdetailed
\fi

\newif\ifwithsols
\ifdefined\setwithsols
  \setwithsols
\fi

\newif\ifwhiteblack
\ifdefined\setwhiteblack
  \setwhiteblack
\fi

% Document-class specific packages (loaded before common_preamble.tex / polyglossia)
\usepackage{fancyhdr}
\usepackage{enumitem}
\usepackage{titlesec}
\usepackage[bottom]{footmisc}
\usepackage{titling}
\usepackage{longtable} % page-breakable tables
\usepackage{float}
\usepackage[onehalfspacing]{setspace}

\ifwhiteblack
  \usepackage[a4paper, landscape, top=1.25cm, bottom=1.25cm, left=1.25cm, right=1.25cm, headheight=1.5cm, headsep=0.5cm]{geometry}
\else
  \usepackage[a4paper, top=2.4cm, bottom=2.4cm, left=2.2cm, right=2.2cm, headheight=1.4cm, headsep=0.4cm, footskip=1.0cm]{geometry}
\fi

% Shared common packages, fonts, polyglossia, callout boxes, code helpers, date & header macros
\IfFileExists{scripts/common_preamble.tex}{%
  % =============================================================================
% Shared Common Preamble (scripts/common_preamble.tex)
% Included by:
%   - scripts/tex_preamble.tex    (worksheets, homeworks, exams, lesson notes)
%   - scripts/book_preamble.tex   (two-sided Arabic book)
%   - scripts/beamer_preamble.tex (Beamer presentation slides)
% =============================================================================

\usepackage{fontspec}
\usepackage{amsmath}
\usepackage{amssymb}
\usepackage{xcolor}
\usepackage{graphicx}
\usepackage{listings}
\usepackage[most]{tcolorbox}
\usepackage{pifont}
\usepackage{minted}
\usepackage{etoolbox}
\usepackage{array}
\usepackage{booktabs}
\usepackage{tikz}
\usepackage{datetime2}

\usetikzlibrary{automata, positioning, arrows.meta, arrows, shapes.geometric, calc}

% Beamer already loads hyperref internally; load it here for article/book classes
% right before polyglossia (bidi).
\makeatletter
\@ifclassloaded{beamer}{}{%
  \usepackage{hyperref}%
  \hypersetup{hidelinks}%
  \providecommand{\LTcaptype}{table}%
}
\makeatother

% Load polyglossia (and bidi) AFTER graphics, tables, boxes, code, and hyperref
\usepackage{polyglossia}
\setmainlanguage[numerals=western]{arabic}
\setotherlanguage{english}
\newfontfamily\arabicfont[Script=Arabic]{Calibri}
\newfontfamily\arabicfontsf[Script=Arabic]{Calibri}
\newfontfamily\arabicfonttt[Script=Arabic]{Courier New}
\newfontfamily\emoji[BoldFont=*,ItalicFont=*]{Segoe UI Emoji}

% Allow slight inter-word stretch to avoid Overfull \hbox with inline \textenglish
\emergencystretch=1.5em

% -----------------------------------------------------------------------------
% Default Minted & Callout Box Definitions
% -----------------------------------------------------------------------------
\AtBeginEnvironment{minted}{\setstretch{0.95}}
\setminted{fontsize=\small, breaklines=true, vspace=2pt}

\tcbset{
  arc=2.5pt, outer arc=2.5pt,
  boxsep=1pt, left=6pt, right=6pt, top=3pt, bottom=3pt,
  toptitle=1.5pt, bottomtitle=1.5pt,
  before skip=6pt, after skip=6pt,
  fonttitle=\bfseries,
  fontupper=\setstretch{1.1}\setlength{\parskip}{2pt},
  colback=white, colframe=black, boxrule=0.6pt,
  boxcolor/.style 2 args={colback=#1, colframe=#2}
}

\newtcolorbox{boxDef}[1][]{boxcolor={blue!4!white}{blue!75!black},
  title={{\emoji📘} تعريف\ifx\\#1\\\else ~#1\fi :}}
\newtcolorbox{boxExercise}[1][]{boxcolor={cyan!4!white}{cyan!70!black},
  title={{\emoji🧩} تمرين\ifx\\#1\\\else ~#1\fi :}}
\newtcolorbox{boxExample}[1][]{boxcolor={yellow!6!white}{orange!85!black},
  title={{\emoji📝} مثال\ifx\\#1\\\else ~#1\fi :}}
\newtcolorbox{boxNote}[1][]{boxcolor={gray!7!white}{black!75},
  title={{\emoji✨} ملاحظة\ifx\\#1\\\else ~#1\fi :}}
\newtcolorbox{boxAttention}[1][]{boxcolor={magenta!5!white}{magenta!75!black},
  title={{\emoji🔔} تنبيه\ifx\\#1\\\else ~#1\fi :}}
\newtcolorbox{boxWarning}[1][]{boxcolor={red!4!white}{red!75!black},
  title={{\emoji⚡} ملاحظة هامة\ifx\\#1\\\else ~#1\fi :}}
\newtcolorbox{boxSolution}[1][]{boxcolor={green!4!white}{green!60!black},
  title={{\emoji✅} حل\ifx\\#1\\\else ~#1\fi :}}
\newtcolorbox{boxSymbol}[1][]{boxcolor={purple!4!white}{purple!70!black},
  title={{\emoji🔣} رمز\ifx\\#1\\\else ~#1\fi :}}
\newtcolorbox{boxHint}[1][]{boxcolor={teal!4!white}{teal!65!black},
  title={{\emoji💡} تلميح\ifx\\#1\\\else ~#1\fi :}}

% Compact code boxes: tighter line spacing and small vertical gaps
\tcbset{
  simplecode/.style={
    arc=2pt, outer arc=2pt,
    left=4pt, right=4pt, top=4pt, bottom=4pt,
    fontupper=\setstretch{0.95},
    before skip=6pt, after skip=6pt,
    colback=gray!5, colframe=black!50, boxrule=0.4pt
  }
}
\newenvironment{boxCode}{\begin{tcolorbox}[simplecode]}{\end{tcolorbox}}

% Code cell used inside tables (keeps code left-to-right inside RTL tables)
% Use inside a fixed-width p{...} column: \begin{pycell}\begin{minted}...\end{minted}\end{pycell}
\makeatletter
\newtcolorbox{pycell}{frame empty,colback=white,boxsep=0pt,left=2pt,right=2pt,top=2pt,bottom=2pt,width=\linewidth,before=\noindent,after=,baseline={\tcb@height-\ht\strutbox},fontupper=\setstretch{0.95}\setlength{\topsep}{0pt}\setlength{\partopsep}{0pt}\setlength{\parskip}{0pt}}
\makeatother

\newcolumntype{C}[1]{>{\centering\arraybackslash}p{#1}}

% Dotted date style (DD.MM.YYYY)
\makeatletter
\DTMnewdatestyle{dotted}{%
  \renewcommand{\DTMdisplaydate}[4]{\two@digits{##3}.\two@digits{##2}.##1}%
  \renewcommand{\DTMDisplaydate}{\DTMdisplaydate}%
}
\makeatother
\DTMsetdatestyle{dotted}

% -----------------------------------------------------------------------------
% Customizable Header / Footer Metadata
% -----------------------------------------------------------------------------
\newcommand{\headerleft}{مدرسة التسامح الشاملة}
\newcommand{\headerright}{الأستاذ محمود اغبارية}
\newcommand{\footercenter}{\thepage}

\newcommand{\setheaderleft}[1]{\renewcommand{\headerleft}{#1}}
\newcommand{\setheaderright}[1]{\renewcommand{\headerright}{#1}}
\newcommand{\setfootercenter}[1]{\renewcommand{\footercenter}{#1}}
%
}{%
  \IfFileExists{../../../scripts/common_preamble.tex}{%
    % =============================================================================
% Shared Common Preamble (scripts/common_preamble.tex)
% Included by:
%   - scripts/tex_preamble.tex    (worksheets, homeworks, exams, lesson notes)
%   - scripts/book_preamble.tex   (two-sided Arabic book)
%   - scripts/beamer_preamble.tex (Beamer presentation slides)
% =============================================================================

\usepackage{fontspec}
\usepackage{amsmath}
\usepackage{amssymb}
\usepackage{xcolor}
\usepackage{graphicx}
\usepackage{listings}
\usepackage[most]{tcolorbox}
\usepackage{pifont}
\usepackage{minted}
\usepackage{etoolbox}
\usepackage{array}
\usepackage{booktabs}
\usepackage{tikz}
\usepackage{datetime2}

\usetikzlibrary{automata, positioning, arrows.meta, arrows, shapes.geometric, calc}

% Beamer already loads hyperref internally; load it here for article/book classes
% right before polyglossia (bidi).
\makeatletter
\@ifclassloaded{beamer}{}{%
  \usepackage{hyperref}%
  \hypersetup{hidelinks}%
  \providecommand{\LTcaptype}{table}%
}
\makeatother

% Load polyglossia (and bidi) AFTER graphics, tables, boxes, code, and hyperref
\usepackage{polyglossia}
\setmainlanguage[numerals=western]{arabic}
\setotherlanguage{english}
\newfontfamily\arabicfont[Script=Arabic]{Calibri}
\newfontfamily\arabicfontsf[Script=Arabic]{Calibri}
\newfontfamily\arabicfonttt[Script=Arabic]{Courier New}
\newfontfamily\emoji[BoldFont=*,ItalicFont=*]{Segoe UI Emoji}

% Allow slight inter-word stretch to avoid Overfull \hbox with inline \textenglish
\emergencystretch=1.5em

% -----------------------------------------------------------------------------
% Default Minted & Callout Box Definitions
% -----------------------------------------------------------------------------
\AtBeginEnvironment{minted}{\setstretch{0.95}}
\setminted{fontsize=\small, breaklines=true, vspace=2pt}

\tcbset{
  arc=2.5pt, outer arc=2.5pt,
  boxsep=1pt, left=6pt, right=6pt, top=3pt, bottom=3pt,
  toptitle=1.5pt, bottomtitle=1.5pt,
  before skip=6pt, after skip=6pt,
  fonttitle=\bfseries,
  fontupper=\setstretch{1.1}\setlength{\parskip}{2pt},
  colback=white, colframe=black, boxrule=0.6pt,
  boxcolor/.style 2 args={colback=#1, colframe=#2}
}

\newtcolorbox{boxDef}[1][]{boxcolor={blue!4!white}{blue!75!black},
  title={{\emoji📘} تعريف\ifx\\#1\\\else ~#1\fi :}}
\newtcolorbox{boxExercise}[1][]{boxcolor={cyan!4!white}{cyan!70!black},
  title={{\emoji🧩} تمرين\ifx\\#1\\\else ~#1\fi :}}
\newtcolorbox{boxExample}[1][]{boxcolor={yellow!6!white}{orange!85!black},
  title={{\emoji📝} مثال\ifx\\#1\\\else ~#1\fi :}}
\newtcolorbox{boxNote}[1][]{boxcolor={gray!7!white}{black!75},
  title={{\emoji✨} ملاحظة\ifx\\#1\\\else ~#1\fi :}}
\newtcolorbox{boxAttention}[1][]{boxcolor={magenta!5!white}{magenta!75!black},
  title={{\emoji🔔} تنبيه\ifx\\#1\\\else ~#1\fi :}}
\newtcolorbox{boxWarning}[1][]{boxcolor={red!4!white}{red!75!black},
  title={{\emoji⚡} ملاحظة هامة\ifx\\#1\\\else ~#1\fi :}}
\newtcolorbox{boxSolution}[1][]{boxcolor={green!4!white}{green!60!black},
  title={{\emoji✅} حل\ifx\\#1\\\else ~#1\fi :}}
\newtcolorbox{boxSymbol}[1][]{boxcolor={purple!4!white}{purple!70!black},
  title={{\emoji🔣} رمز\ifx\\#1\\\else ~#1\fi :}}
\newtcolorbox{boxHint}[1][]{boxcolor={teal!4!white}{teal!65!black},
  title={{\emoji💡} تلميح\ifx\\#1\\\else ~#1\fi :}}

% Compact code boxes: tighter line spacing and small vertical gaps
\tcbset{
  simplecode/.style={
    arc=2pt, outer arc=2pt,
    left=4pt, right=4pt, top=4pt, bottom=4pt,
    fontupper=\setstretch{0.95},
    before skip=6pt, after skip=6pt,
    colback=gray!5, colframe=black!50, boxrule=0.4pt
  }
}
\newenvironment{boxCode}{\begin{tcolorbox}[simplecode]}{\end{tcolorbox}}

% Code cell used inside tables (keeps code left-to-right inside RTL tables)
% Use inside a fixed-width p{...} column: \begin{pycell}\begin{minted}...\end{minted}\end{pycell}
\makeatletter
\newtcolorbox{pycell}{frame empty,colback=white,boxsep=0pt,left=2pt,right=2pt,top=2pt,bottom=2pt,width=\linewidth,before=\noindent,after=,baseline={\tcb@height-\ht\strutbox},fontupper=\setstretch{0.95}\setlength{\topsep}{0pt}\setlength{\partopsep}{0pt}\setlength{\parskip}{0pt}}
\makeatother

\newcolumntype{C}[1]{>{\centering\arraybackslash}p{#1}}

% Dotted date style (DD.MM.YYYY)
\makeatletter
\DTMnewdatestyle{dotted}{%
  \renewcommand{\DTMdisplaydate}[4]{\two@digits{##3}.\two@digits{##2}.##1}%
  \renewcommand{\DTMDisplaydate}{\DTMdisplaydate}%
}
\makeatother
\DTMsetdatestyle{dotted}

% -----------------------------------------------------------------------------
% Customizable Header / Footer Metadata
% -----------------------------------------------------------------------------
\newcommand{\headerleft}{مدرسة التسامح الشاملة}
\newcommand{\headerright}{الأستاذ محمود اغبارية}
\newcommand{\footercenter}{\thepage}

\newcommand{\setheaderleft}[1]{\renewcommand{\headerleft}{#1}}
\newcommand{\setheaderright}[1]{\renewcommand{\headerright}{#1}}
\newcommand{\setfootercenter}[1]{\renewcommand{\footercenter}{#1}}
%
  }{%
    \input{../../scripts/common_preamble.tex}%
  }%
}

\pagestyle{fancy}

% Prevent widow/orphan lines
\widowpenalty=10000
\clubpenalty=10000

% Compact section headings and default list spacing so pages do not waste vertical space
\titlespacing*{\section}{0pt}{1.4ex plus 0.4ex minus 0.2ex}{0.7ex plus 0.2ex}
\titlespacing*{\subsection}{0pt}{1.1ex plus 0.3ex minus 0.2ex}{0.5ex plus 0.1ex}
\titlespacing*{\subsubsection}{0pt}{0.9ex plus 0.2ex minus 0.1ex}{0.4ex plus 0.1ex}
\setlist{topsep=4pt, partopsep=0pt, parsep=2pt, itemsep=4pt}

\gappto\captionsarabic{\renewcommand{\contentsname}{جدول المحتويات}}

% Legacy listings settings
\lstset{
  language=[Sharp]C,
  numbers=left,
  stepnumber=1,
  numbersep=8pt,
  frame=single,
  basicstyle=\ttfamily\small
}
\ifwhiteblack
  \lstset{
    keywordstyle=\bfseries,
    stringstyle=\ttfamily,
    commentstyle=\itshape
  }
\else
  \lstset{
    keywordstyle=\color{blue},
    stringstyle=\color{red},
    commentstyle=\color{green!50!black}
  }
\fi

% White-and-black landscape presentation mode overrides
\ifwhiteblack
  \AtBeginDocument{
    \LARGE
    \usemintedstyle{bw}
  }
  \setminted{vspace=2pt}
  \tcbset{
    colback=white, colframe=black, boxrule=1pt,
    boxcolor/.style 2 args={colback=white, colframe=black},
    simplecode/.append style={colback=white, colframe=black, boxrule=1pt}
  }
\fi

% Default to empty author to avoid "LaTeX Warning: No \author given" when omitted
\author{}

% Compact title block
\makeatletter
\renewcommand{\@maketitle}{%
  \begin{center}%
    {\ifwhiteblack\Huge\else\LARGE\fi \bfseries \@title \par}%
    \ifdefempty{\@author}{}{%
      \vskip 0.15em%
      {\large \@author \par}%
    }%
    {\small \hfill \textenglish{\DTMtoday} \par}%
  \end{center}%
  \vspace{-0.3cm}%
}
\makeatother

\apptocmd{\maketitle}{\thispagestyle{fancy}}{}{}

% Apply page headers and footer (using \headerleft, \headerright, \footercenter from common_preamble.tex)
\newcommand{\applyfancyheaders}{%
  \fancyhf{}%
  \fancyhead[L]{\textarabic{\headerleft}}%
  \fancyhead[R]{\textarabic{\headerright}}%
  \fancyfoot[C]{\textarabic{\footercenter}}%
}

\fancypagestyle{plain}{\applyfancyheaders}
\applyfancyheaders

\setcounter{tocdepth}{3} % only section, subsection, and subsubsection in TOC

\newcommand{\importpdfpage}[2]{%
    \newgeometry{left=0cm,right=0cm,top=0cm,bottom=0cm}
    \thispagestyle{empty}
    \includegraphics[page=#2,width=0.95\paperwidth,height=0.95\paperheight,keepaspectratio]{#1}
    \restoregeometry
}

\newcommand{\insertFullImg}[1]{%
    \noindent
    \makebox[\textwidth][c]{
        \includegraphics[width=0.9\paperwidth,keepaspectratio]{#1}%
    }%
}%


\ifwithsols
\title{حل امتحان تصنيف في علم الحاسوب \\ الصف التاسع د}
\else
\title{امتحان تصنيف في علم الحاسوب \\ الصف التاسع د}
\fi
\renewcommand{\DTMtoday}{12.11.2026}

\begin{document}

\maketitle
\thispagestyle{fancy}

\ifwithsols
\else
\begin{boxCode}
    \begin{itemize}
        \item يحتوي الاختبار على 6 أسئلة، عليك حل جميع الأسئلة.
        \item الوقت المخصص: ساعة ونصف.
        \item يسمح باستخدام كل مادة مساعدة، عدا الآلة التي يمكن برمجتها.
        \item اكتب بقلم حبر فقط، وبخط واضح.
        \item في كل الأسئلة افترض أن المدخلات صحيحة ولا داعي لفحصها إلا إذا ذُكر غير ذلك.
    \end{itemize}
\end{boxCode}

\vspace{1cm}
\textbf{الاسم:} \rule{5cm}{0.4pt}

\vspace{1cm}
\begin{center}
  \includegraphics[width=0.4\textwidth]{../../../images/logoTasamoh.png}
\end{center}

\vspace{1cm}
\begin{flushleft}
    نرجو لكم النجاح والتفوق.

    معلمو الموضوع.
\end{flushleft}

\clearpage
\fi

\begin{enumerate}[label=\textbf{السؤال \arabic*:}, wide=0pt, itemsep=2.2em]

% ---------------------------------------------------------------
\item ($16$ علامة، $4$ علامات لكل بند) أمامكم جدول فيه أربعة مقاطع برنامج بلغة بايثون. لكل مقطع برنامج في الجدول، اكتبوا ما هو المُخرج الناتج في العمود الفارغ.

{\renewcommand{\arraystretch}{1.4}
\begin{longtable}{|C{1cm}|p{8cm}|C{5cm}|}
\hline
\textbf{البند} & \centering\textbf{كود بلغة بايثون} & \textbf{المُخرج الناتج} \\
\hline
$1$ &
\begin{pycell}
\begin{english}
\begin{minted}{python}
a = 7
b = a + 3
a = b - 5
print("a = ", a)
print("b = ", b)
\end{minted}
\end{english}
\end{pycell}
& \ifwithsols \textenglish{\small\texttt{a = 5}}\newline\textenglish{\small\texttt{b = 10}} \fi \\
\hline
$2$ &
\begin{pycell}
\begin{english}
\begin{minted}{python}
value = 15
result = value // 2
value = result + 3
print("value = ", value, end=" ")
print("result = ", result)
\end{minted}
\end{english}
\end{pycell}
& \ifwithsols \textenglish{\small\texttt{value = 10 result = 7}} \fi \\
\hline
$3$ &
\begin{pycell}
\begin{english}
\begin{minted}{python}
num = 7890
num1 = (num // 100) % 10
num2 = num % 100
print("num1 = ", num1, end=",")
print("num2 = ", num2)
\end{minted}
\end{english}
\end{pycell}
& \ifwithsols \textenglish{\small\texttt{num1 = 8,num2 = 90}} \fi \\
\hline
$4$ &
\begin{pycell}
\begin{english}
\begin{minted}{python}
num = 47
x = num // 10 + 10 * (num % 10)
print("num = ", num)
print("x = ", x)
\end{minted}
\end{english}
\end{pycell}
& \ifwithsols \textenglish{\small\texttt{num = 47}}\newline\textenglish{\small\texttt{x = 74}} \fi \\
\hline
\end{longtable}}

\clearpage

% ---------------------------------------------------------------
\item ($16$ علامة، $4$ علامات لكل بند) في الجدول الذي أمامكم أربعة مقاطع برنامج جزئية والمُخرج الذي سينتج من تشغيل مقطع البرنامج الكامل. كل مقطع برنامج يحتوي على حلقة.
أكملوا التعليمات الناقصة في كل مقطع برنامج للحصول على المُخرج المعطى. استعينوا بالمثال في أعلى الجدول.

{\renewcommand{\arraystretch}{1.4}
\begin{longtable}{|C{1.5cm}|p{9cm}|C{4cm}|}
\hline
\textbf{البند} & \centering\textbf{كود بلغة بايثون} & \textbf{المُخرج الناتج} \\
\hline
مثال &
\begin{pycell}
\begin{english}
\begin{minted}{python}
for item in range(2, 7, 2):
    print(item, end=" ")
\end{minted}
\end{english}
\end{pycell}
& \textenglish{\texttt{2 4 6}} \\
\hline
$1$ &
\begin{pycell}
\begin{english}
\begin{minted}{python}
for item in range(2, ________):
    print(item, end=" ")
\end{minted}
\end{english}
\end{pycell}
& \textenglish{\texttt{2 3}} \\
\hline
$2$ &
\begin{pycell}
\begin{english}
\begin{minted}{python}
for item in range(10):
    if item % 3 == ________________:
        print(item, end=" ")
\end{minted}
\end{english}
\end{pycell}
& \textenglish{\texttt{1 4 7}} \\
\hline
$3$ &
\begin{pycell}
\begin{english}
\begin{minted}{python}
item = _________________
while item < 10:
    print(item, end=" ")
    item = __________________
\end{minted}
\end{english}
\end{pycell}
& \textenglish{\texttt{2 5 8}} \\
\hline
$4$ &
\begin{pycell}
\begin{english}
\begin{minted}{python}
item = 3
while __________ > __________:
    print(item * 5)
    item = item - 1
\end{minted}
\end{english}
\end{pycell}
& \textenglish{\texttt{15}}\newline\textenglish{\texttt{10}}\newline\textenglish{\texttt{5}} \\
\hline
\end{longtable}}

\ifwithsols
\begin{boxSolution}
\begin{enumerate}[label=\arabic*.]
\item \textenglish{\texttt{for item in range(2, 4):}}
\item \textenglish{\texttt{if item \% 3 == 1:}}
\item التهيئة \textenglish{\texttt{item = 2}}، والتحديث داخل الحلقة \textenglish{\texttt{item = item + 3}}
\item الشرط \textenglish{\texttt{while item > 0:}}
\end{enumerate}
\end{boxSolution}
\fi
\filbreak

% ---------------------------------------------------------------
\item ($16$ علامة) أمامكم عمودان من مقاطع البرنامج التي تُنفَّذ فيها حلقة. يجب الملاءمة بين أزواج الحلقات المتكافئة (بحيث ينتج في مقطعي الحلقة نفس المُخرج تمامًا) وكتابتها في الجدول أدناه.

{\renewcommand{\arraystretch}{1.3}
\begin{longtable}{|C{0.9cm}|p{5.8cm}|C{0.9cm}|p{6cm}|}
\hline
\multicolumn{2}{|c|}{\textbf{العمود أ}} & \multicolumn{2}{c|}{\textbf{العمود ب}} \\
\hline
$1$ &
\begin{pycell}
\begin{english}
\begin{minted}{python}
i = 5
while i <= 30:
    print(i)
    i = i + 5
\end{minted}
\end{english}
\end{pycell}
& $A$ &
\begin{pycell}
\begin{english}
\begin{minted}{python}
for i in range(16, 2, -3):
    print(i)
\end{minted}
\end{english}
\end{pycell}
\\ \hline
$2$ &
\begin{pycell}
\begin{english}
\begin{minted}{python}
for i in range(5, 10):
    print(i)
\end{minted}
\end{english}
\end{pycell}
& $B$ &
\begin{pycell}
\begin{english}
\begin{minted}{python}
x = 2
while x < 13:
    print(x)
    x = x + 2
\end{minted}
\end{english}
\end{pycell}
\\ \hline
$3$ &
\begin{pycell}
\begin{english}
\begin{minted}{python}
i = 1
while i <= 6:
    print(i * 2)
    i = i + 1
\end{minted}
\end{english}
\end{pycell}
& $C$ &
\begin{pycell}
\begin{english}
\begin{minted}{python}
for i in range(5, 31, 5):
    print(i)
\end{minted}
\end{english}
\end{pycell}
\\ \hline
$4$ &
\begin{pycell}
\begin{english}
\begin{minted}{python}
i = 16
while i > 2:
    print(i)
    i = i - 3
\end{minted}
\end{english}
\end{pycell}
& $D$ &
\begin{pycell}
\begin{english}
\begin{minted}{python}
item = 4
while item < 9:
    print(item + 1)
    item = item + 1
\end{minted}
\end{english}
\end{pycell}
\\ \hline
\end{longtable}}

{\renewcommand{\arraystretch}{1.3}
\begin{longtable}{|C{4.5cm}|C{5cm}|}
\hline
\textbf{رقم الحلقة في العمود أ} & \textbf{الحلقة المكافئة من العمود ب ($A$ / $B$ / $C$ / $D$)} \\
\hline
$1$ & \ifwithsols $\mathbf{C}$ (مضاعفات الـ $5$ من $5$ إلى $30$) \fi \\ \hline
$2$ & \ifwithsols $\mathbf{D}$ (الأعداد من $5$ إلى $9$) \fi \\ \hline
$3$ & \ifwithsols $\mathbf{B}$ (الأعداد الزوجية من $2$ إلى $12$) \fi \\ \hline
$4$ & \ifwithsols $\mathbf{A}$ (الأعداد $16, 13, 10, 7, 4$) \fi \\ \hline
\end{longtable}}
\filbreak

% ---------------------------------------------------------------
\item ($16$ علامة) في لعبة \textbf{"التنين الذهبي"}، يحصل اللاعب على رتبة قوّة وفقًا للنقاط التي جمعها في اللعبة. رتب القوّة الممكنة هي:
محارب مبتدئ (\textenglish{Beginner Warrior})، فارس فضّي (\textenglish{Silver Knight})، فارس ذهبي (\textenglish{Gold Knight})، وساحر خارق (\textenglish{Super Wizard}).
اللاعب الذي لديه أقل من $50$ نقطة يُعرَّف برتبة "مبتدئ" (\textenglish{Beginner}).
تُحدَّد رتبة القوّة وفقًا للسيناريو الذي أمامكم:

\begin{boxCode}
\begin{english}
\begin{minted}{python}
points = int(input("Enter Point: "))
if points < 50:
    print("Beginner")
elif points < 200:
    print("Beginner Warrior")
elif points < 500:
    print("Silver Knight")
elif points < 1000:
    print("Gold Knight")
else:
    print("Super Wizard")
\end{minted}
\end{english}
\end{boxCode}

\begin{enumerate}[label=\alph*.]
\item ($6$ علامات) معطى الجدول التالي. اكتبوا ما هي رتبة القوّة للاعب لكل واحدة من قيم النقاط في الجدول:

{\renewcommand{\arraystretch}{1.4}
\begin{longtable}{|C{3cm}|C{6cm}|}
\hline
\textbf{النقاط} & \textbf{رتبة القوّة} \\
\hline
$25$ & \ifwithsols مبتدئ (\textenglish{Beginner}) \fi \\ \hline
$120$ & \ifwithsols محارب مبتدئ (\textenglish{Beginner Warrior}) \fi \\ \hline
$200$ & \ifwithsols فارس فضّي (\textenglish{Silver Knight}) \fi \\ \hline
$500$ & \ifwithsols فارس ذهبي (\textenglish{Gold Knight}) \fi \\ \hline
$555$ & \ifwithsols فارس ذهبي (\textenglish{Gold Knight}) \fi \\ \hline
$1050$ & \ifwithsols ساحر خارق (\textenglish{Super Wizard}) \fi \\ \hline
\end{longtable}}
\filbreak

\item ($10$ علامات) معطى الجدول التالي. أكملوا "رتبة القوّة" وفقًا للبرنامج (رتب القوّة الممكنة: مبتدئ، محارب مبتدئ، فارس فضّي، فارس ذهبي، ساحر خارق):

{\renewcommand{\arraystretch}{1.4}
\begin{longtable}{|C{1cm}|>{\raggedright\arraybackslash}p{9cm}|C{4.5cm}|}
\hline
\textbf{البند} & \centering\textbf{مجال النقاط} & \textbf{رتبة القوّة} \\
\hline
مثال & \textenglish{\texttt{points == 10}} & مبتدئ (\textenglish{Beginner}) \\ \hline
$1$ & \textenglish{\texttt{points < 1000 and points >= 500}} & \ifwithsols فارس ذهبي (\textenglish{Gold Knight}) \fi \\ \hline
$2$ & \textenglish{\texttt{points >= 200 and points < 500}} & \ifwithsols فارس فضّي (\textenglish{Silver Knight}) \fi \\ \hline
$3$ & \textenglish{\texttt{points < 200 and points >= 50}} & \ifwithsols محارب مبتدئ (\textenglish{Beginner Warrior}) \fi \\ \hline
$4$ & \textenglish{\texttt{points >= 1000}} & \ifwithsols ساحر خارق (\textenglish{Super Wizard}) \fi \\ \hline
$5$ & \textenglish{\texttt{points < 50 and points >= 0}} & \ifwithsols مبتدئ (\textenglish{Beginner}) \fi \\ \hline
\end{longtable}}
\filbreak
\end{enumerate}
\filbreak

% ---------------------------------------------------------------
\item ($15$ علامة) في لعبة \textbf{"صائد الكنوز"}، تُعتبر النقاط "جيدة" إذا كانت \textbf{فوق $100$ وحتى $300$ (يشمل $300$)}.
ينفّذ اللاعب $10$ مهام مختلفة، وفي كل مهمة يحصل على نقاط وفقًا لنجاحه.
يجب على اللعبة أن تعدّ وتطبع كم مرة حقّق اللاعب نقاطًا "جيدة".
أمامكم ثلاثة مقاطع برنامج. حدّدوا لكل مقطع برنامج هل ينفّذ المطلوب في السؤال. إذا لا، اشرحوا كيف يمكن تصحيح الكود لكي ينفّذ المطلوب.

\begin{enumerate}[label=\textbf{مقطع \alph*:}, wide=0pt, itemsep=1.5em]

\item \mbox{}
\begin{boxCode}
\begin{english}
\begin{minted}{python}
count = 0
for i in range(10):
    score = int(input("Enter score: "))
    if score > 100 and score <= 300:
        count = count + 1
print(count)
\end{minted}
\end{english}
\end{boxCode}
\ifwithsols
\begin{boxSolution}
\textbf{نعم، ينفّذ المطلوب.} الحلقة تتكرّر $10$ مرات، تستقبل النقاط في كل مهمة، وتفحص بشكل صحيح إذا كانت أكبر من $100$ وأصغر من أو تساوي $300$.
\end{boxSolution}
\else
\vspace{0.5cm}
\textbf{هل ينفّذ الكود المطلوب؟} \hrulefill \\
 \\
\textbf{طريقة التصحيح (إن لزم):} \hrulefill
\fi
\filbreak

\item \mbox{}
\begin{boxCode}
\begin{english}
\begin{minted}{python}
count = 0
for i in range(10):
    score = int(input("Enter score: "))
    if score > 100 or score <= 300:
        count = count + 1
print(count)
\end{minted}
\end{english}
\end{boxCode}
\ifwithsols
\begin{boxSolution}
\textbf{لا ينفّذ المطلوب.} لأنّه استعمل العامل \textenglish{or} بدل \textenglish{and}، فيتحقّق الشرط لكل عدد مهما كانت قيمته. \\
\textbf{التصحيح:} تغيير \textenglish{or} إلى \textenglish{and} في سطر الشرط: \textenglish{\texttt{if score > 100 and score <= 300:}}.
\end{boxSolution}
\else
\vspace{0.5cm}
\textbf{هل ينفّذ الكود المطلوب؟} \hrulefill \\
 \\
\textbf{طريقة التصحيح (إن لزم):} \hrulefill
\fi
\filbreak

\item \mbox{}
\begin{boxCode}
\begin{english}
\begin{minted}{python}
count = 0
for i in range(1, 10, 1):
    score = int(input("Enter score: "))
    if score > 100:
        if score <= 300:
            count = count + 1
print(count)
\end{minted}
\end{english}
\end{boxCode}
\ifwithsols
\begin{boxSolution}
\textbf{لا ينفّذ المطلوب.} لأنّ الحلقة \textenglish{\texttt{range(1, 10, 1)}} تتكرّر $9$ مرات فقط (من $1$ إلى $9$) بدل $10$ مرات. \\
\textbf{التصحيح:} تغيير مدى الحلقة إلى \textenglish{\texttt{range(10)}} أو \textenglish{\texttt{range(1, 11, 1)}}.
\end{boxSolution}
\else
\vspace{0.5cm}
\textbf{هل ينفّذ الكود المطلوب؟} \hrulefill \\
 \\
\textbf{طريقة التصحيح (إن لزم):} \hrulefill
\fi
\filbreak

\end{enumerate}

\clearpage

% ---------------------------------------------------------------
% \item ($20$ علامة) أمامكم مقطع برنامج يستقبل عددين صحيحين موجبين وينفّذ حسابًا معيّنًا:

% \begin{boxCode}
% \begin{english}
% \begin{minted}{python}
% num1 = int(input('Please enter number '))
% num2 = int(input('Please enter number '))
% while num1 >= num2:
%     num1 = num1 - num2
% print(num1)
% \end{minted}
% \end{english}
% \end{boxCode}

% \begin{enumerate}[label=\alph*.]
% \item أكملوا الجدول التالي لقيم \textenglish{num1} و \textenglish{num2} المعطاة:

% {\begin{english}
% \renewcommand{\arraystretch}{1.4}
% \begin{longtable}{|C{3cm}|C{3cm}|C{3.5cm}|}
% \hline
% \textbf{num1} & \textbf{num2} & \textarabic{\textbf{المُخرج}} \\
% \hline
% $11$ & $4$ & \ifwithsols $\mathbf{3}$ \fi \\ \hline
% $9$ & $3$ & \ifwithsols $\mathbf{0}$ \fi \\ \hline
% $5$ & $2$ & \ifwithsols $\mathbf{1}$ \fi \\ \hline
% $2$ & $5$ & \ifwithsols $\mathbf{2}$ \fi \\ \hline
% \end{longtable}
% \end{english}}
% \filbreak

% \item ماذا ينفّذ مقطع البرنامج؟ أحيطوا بدائرة الإجابة الصحيحة:
% \begin{enumerate}[label=\arabic*.]
% \item يحسب ويطبع حاصل ضرب \textenglish{num1} بـ \textenglish{num2}.
% \item يحسب ويطبع الفرق بين \textenglish{num1} و \textenglish{num2}.
% \item يحسب ويطبع باقي القسمة (المودولو) من قسمة \textenglish{num1} على \textenglish{num2}.
% \item يحسب ويطبع \textenglish{num1} مرفوع للقوة \textenglish{num2}.
% \end{enumerate}

% \ifwithsols
% \begin{boxSolution}
% الإجابة الصحيحة هي \textbf{$3$}: يحسب ويطبع باقي القسمة (المودولو) من قسمة \textenglish{num1} على \textenglish{num2} عن طريق الطرح المتكرّر لـ \textenglish{num2} من \textenglish{num1} طالما أنّ $\text{num1} \geq \text{num2}$.
% \end{boxSolution}
% \fi
% \end{enumerate}

% \filbreak

% ---------------------------------------------------------------
\item ($21$ علامة) \textbf{العدد "الموقِظ"} هو عدد صحيح موجب يكون فيه مجموع جميع منازله \textbf{أكبر بـ $3$ أضعاف أو أكثر} من عدد المنازل في العدد.
\begin{itemize}
\item \textbf{مثال $1$:} العدد $129$ هو عدد "موقِظ" لأنّ مجموع منازله ($1 + 2 + 9 = 12$) أكبر بـ $3$ أضعاف أو أكثر من عدد منازله ($3$). الحساب: $12 / 3 \geq 3$.
\item \textbf{مثال $2$:} العدد $5$ هو عدد "موقِظ" لأنّ مجموع منازله ($5$) أكبر بـ $3$ أضعاف أو أكثر من عدد منازله ($1$). الحساب: $5 / 1 \geq 3$.
\item \textbf{مثال $3$:} العدد $124$ ليس عددًا "موقِظًا" لأنّ مجموع منازله ($1 + 2 + 4 = 7$) ليس أكبر بـ $3$ أضعاف أو أكثر من عدد منازله ($3$). الحساب: $7 / 3 < 3$.
\end{itemize}

\begin{enumerate}[label=\alph*., itemsep=1.5em]
\item ($6$ علامات — $3$ علامات لكل بند) حسب التعريف أعلاه، حدّدوا لكل واحد من العددين التاليين هل هو عدد "موقِظ". إذا نعم، اعرضوا حسابًا رياضيًا (كما في المثال) يُظهر ذلك؛ وإلّا اشرحوا:
\begin{enumerate}[label=\arabic*., itemsep=1em]
\item العدد $4158$:
\ifwithsols
\begin{boxSolution}
\textbf{نعم}، العدد $4158$ هو عدد "موقِظ" لأنّ مجموع منازله هو $4 + 1 + 5 + 8 = 18$، وعدد منازله هو $4$، ويتحقّق:
$$18 / 4 = 4.5 \geq 3 \quad (\text{أو } 18 \geq 3 \cdot 4 = 12)$$
\end{boxSolution}
\else
\hrulefill
\fi

\item العدد $3124$:
\ifwithsols
\begin{boxSolution}
\textbf{لا}، العدد $3124$ ليس عددًا "موقِظًا" لأنّ مجموع منازله هو $3 + 1 + 2 + 4 = 10$، وعدد منازله هو $4$، ولا يتحقّق الشرط:
$$10 / 4 = 2.5 < 3 \quad (\text{أو } 10 < 3 \cdot 4 = 12)$$
\end{boxSolution}
\else
\hrulefill
\fi
\end{enumerate}
\filbreak

\item ($15$ علامة) أمامكم مقطع برنامج يستقبل عددًا من المستخدم ويفحص هل العدد "موقِظ". أكملوا الفراغات $(1)$–$(5)$ في المقطع:

\begin{boxCode}
\begin{english}
\begin{minted}{python}
sum = (1)______________
count = 0
num = int(input("Please enter a number:"))
while num > 0:
    sum = (2)________________ + num % 10
    count = count + (3)______________
    num = num // 10
if (4)_____________ >= 3 * (5)_________________:
    print("The number is awakening")
else:
    print("The number is not awakening")
\end{minted}
\end{english}
\end{boxCode}

\ifwithsols
\begin{boxSolution}
\begin{enumerate}[label=(\arabic*)]
\item $0$
\item \textenglish{\texttt{sum}}
\item $1$
\item \textenglish{\texttt{sum}}
\item \textenglish{\texttt{count}}
\end{enumerate}
\end{boxSolution}
\fi
\filbreak
\end{enumerate}
\filbreak

\end{enumerate}

\end{document}
